\clearpage
\chapter{Introduction}
\label{chap:introduction}

In life and industri that are alot of machinery that need regualted and precice controll, which i why a controller is made to achive each goal as needed. In robotics controllers can be expicaly usefull when trying to complete difficult taskes with multipul moveing parts. Where balance is one of main hindrences, when trying to design such a system. This projeckt will focus trying to solve a balance isuee of an inverted pendulem. 

\section{Project Goals}

With a PLC, motor, conveyor track, inverted pendulem on a cart and a two diffrencet controller. it is attempted to make the inverted pendulem balance rightsideup using a controller to controll the force of the motor. success is mesurd by how well the controller balences the pendulem compared to a simulation of the controller, and the best controller is the one whose closest to the simulation of said controller.

   
\section{Problem Definition}



\section{Constraints \& Limitations}
The limitations for this projeckt come from the physical setup, and the PLC. The cart is limited to move within two sensors on the conveyor track, and the maximum force for the moter is X. For the PLC it cant run faster than 1000Hz, and sampling via OPCUA is limited to 100Hz



\section{System OverView}
